\documentclass[UTF8,fontset=none,11pt]{ctexart}

% ---------- 基础宏包 ----------
\usepackage[a4paper,margin=2.15cm]{geometry}
\usepackage{amsmath,amssymb}
\usepackage{enumitem}
\usepackage{xcolor}
\usepackage{booktabs}
\usepackage{array}
\usepackage{longtable}
\usepackage{tikz}
\usepackage{hyperref}
\usetikzlibrary{positioning,arrows.meta,shapes.geometric,fit,backgrounds,calc}

\setCJKmainfont{PingFang SC}
\setCJKsansfont{PingFang SC}
\setCJKmonofont{PingFang SC}

% ---------- 颜色 ----------
\definecolor{maincol}{RGB}{28,85,150}
\definecolor{accent}{RGB}{185,75,30}
\definecolor{softbg}{RGB}{244,248,252}
\definecolor{defrule}{RGB}{55,135,90}
\definecolor{warnrule}{RGB}{195,130,35}
\definecolor{graytext}{RGB}{105,105,105}

% ---------- 超链接 ----------
\hypersetup{
  colorlinks=true,
  linkcolor=maincol,
  urlcolor=accent,
  pdftitle={语法分析导论、Parser 与 CFG 详细中文学习笔记},
  pdfauthor={Summary}
}

% ---------- 列表风格 ----------
\setlist[itemize]{leftmargin=1.4em,itemsep=2pt,topsep=2pt}
\setlist[enumerate]{leftmargin=1.6em,itemsep=2pt,topsep=2pt}

% ---------- 自定义可分页标注盒子 ----------
\newenvironment{defbox}[1][]{%
  \par\medskip\noindent{\color{defrule}\rule{\linewidth}{1.1pt}}\par\nopagebreak
  \noindent{\bfseries\color{defrule}#1}\par\nopagebreak\smallskip}%
  {\par\nopagebreak\smallskip\noindent{\color{defrule}\rule{\linewidth}{0.4pt}}\par\medskip}

\newenvironment{notebox}[1][]{%
  \par\medskip\noindent{\color{maincol}\rule{\linewidth}{1.1pt}}\par\nopagebreak
  \noindent{\bfseries\color{maincol}#1}\par\nopagebreak\smallskip}%
  {\par\nopagebreak\smallskip\noindent{\color{maincol}\rule{\linewidth}{0.4pt}}\par\medskip}

\newenvironment{warnbox}[1][]{%
  \par\medskip\noindent{\color{warnrule}\rule{\linewidth}{1.1pt}}\par\nopagebreak
  \noindent{\bfseries\color{warnrule}#1}\par\nopagebreak\smallskip}%
  {\par\nopagebreak\smallskip\noindent{\color{warnrule}\rule{\linewidth}{0.4pt}}\par\medskip}

\newcommand{\pg}[1]{\texorpdfstring{{\small\color{graytext}\,[p#1]}}{ [p#1]}}
\newcommand{\term}[2]{\textbf{#1}（#2）}
\newcommand{\code}[1]{\texttt{#1}}

\title{\bfseries 编译原理 · 语法分析导论、Parser 与 CFG\\[2pt]
\large Syntax Analysis: Intro \& Parser \& CFG —— 详细中文学习笔记}
\author{基于课件 \texttt{3.Syntax Analysis\_Intro\_Parser\_CFG.pdf}（中山大学 · 赵帅）整理}
\date{}

\begin{document}
\maketitle
\thispagestyle{empty}

\begin{notebox}[本笔记说明]
本笔记覆盖课件第 3 讲共 55 页内容，主线是：
\[
\text{token 流} \xrightarrow{\text{语法分析 Parser}}
\text{语法结构 Parse Tree / AST}
\]
先说明为什么正则语言不足以描述程序语法，再系统学习\term{文法}{Grammar}、
\term{上下文无关文法}{Context-Free Grammar, CFG}、\term{推导}{Derivation}、
\term{分析树}{Parse Tree}、\term{二义性}{Ambiguity}和
\term{乔姆斯基文法体系}{Chomsky Grammar Hierarchy}。每个重要知识点都标注课件页码 \pg{页码}，
并补充通俗解释与复习用例。
\end{notebox}

\tableofcontents
\vspace{4pt}\hrule

% ============================================================
\section{从词法分析走向语法分析（第 2--5 页）}
% ============================================================

\subsection{前置回顾：词法分析的实现链路 \pg{2,5}}
上一讲已经完成了从\textbf{词法规范}到\textbf{词法分析器实现}的主线：

\begin{center}
\begin{tikzpicture}[font=\small,
  n/.style={draw,rounded corners=2pt,fill=softbg,minimum height=0.72cm,align=center,inner sep=4pt},
  hi/.style={draw,rounded corners=2pt,fill=white,minimum height=0.72cm,align=center,inner sep=4pt},
  arr/.style={-{Latex[length=2mm]},thick}]
\node[n] (re) {正则表达式\\RE};
\node[n,right=1.5cm of re] (nfa) {NFA};
\node[n,right=1.5cm of nfa] (dfa) {DFA};
\node[n,right=1.5cm of dfa] (tab) {表驱动实现\\Table-driven};
\node[hi,below=0.9cm of tab] (lex) {词法分析器\\Lexical Analyzer};
\draw[arr] (re) -- node[above,font=\footnotesize]{Thompson} (nfa);
\draw[arr] (nfa) -- node[above,font=\footnotesize]{子集构造} (dfa);
\draw[arr] (dfa) -- node[above,font=\footnotesize]{最小化} (tab);
\draw[arr] (tab) -- (lex);
\end{tikzpicture}
\end{center}

\begin{itemize}
  \item \term{Thompson 算法}{Thompson Algorithm}：把正则表达式转换为 NFA。
  \item \term{子集构造法}{Subset-Construction Algorithm}：把 NFA 转换为 DFA。
  \item \term{分割法}{Partition Algorithm}：最小化 DFA，减少状态数。
  \item 词法分析的输出是 token 流；本讲从这里继续，研究 token 流是否组成合法程序结构。
\end{itemize}

\subsection{正则语言的局限 \pg{3,4}}
\begin{defbox}[正则语言（Regular Language）的能力边界]
正则语言适合描述\textbf{局部、线性、无需记忆}的模式，例如标识符、数字、字符串、注释等 token。
但它不能表达需要\textbf{任意深度记忆或配对计数}的结构。
\end{defbox}

课件强调了几个典型例子：
\begin{itemize}
  \item \(L=\{a^n b^n \mid n\ge 1\}\) 不是正则语言：识别它需要记住前面读了多少个 \(a\)，再检查后面是否有同样数量的 \(b\)。有限自动机状态数有限，不能无限计数。
  \item \(L=\{a^n b a^n \mid n\ge 0\}\) 也体现同类困难：中间的 \(b\) 左右两侧 \(a\) 的数量要相等。
  \item 括号匹配、花括号匹配 \(\{\{\{\cdots\}\}\}\) 不是正则语言，因为括号可以任意嵌套，必须记住尚未匹配的左括号数量。
  \item 嵌套结构，例如 \code{if ... if ... else ... else}、C 语言中的嵌套 \code{for} 循环，通常不能只靠正则语言描述。
\end{itemize}

\begin{notebox}[直观理解]
有限自动机像一个只有有限记忆的小机器。它可以判断“当前是不是正在读标识符”，但无法记住“已经读过 137 个左括号，所以后面必须出现 137 个右括号”。因此，词法分析可以用正则语言；语法分析需要更强的形式化工具。
\end{notebox}

\subsection{为什么需要语法分析 \pg{4,5}}
程序设计语言的 token 本身通常可以由正则表达式描述；但是 token 之间如何组合成表达式、语句、函数、块结构，则需要能够描述\textbf{递归结构}和\textbf{嵌套结构}的模型。

\[
\text{Regular Language for Lexical Analysis}
\qquad
\text{Context-Free Language for Syntax Analysis}
\]

这就是本讲引入\term{上下文无关文法}{Context-Free Grammar, CFG}的核心原因。

% ============================================================
\section{语法分析在编译器中的位置（第 6--11 页）}
% ============================================================

\subsection{编译阶段中的 Parser \pg{6}}
语法分析是编译器前端的第二个主要阶段，位于词法分析之后、语义分析之前。

\begin{center}
\begin{tikzpicture}[font=\small,
  n/.style={draw,rounded corners=2pt,fill=softbg,minimum height=0.72cm,align=center,inner sep=4pt},
  hi/.style={draw,rounded corners=2pt,fill=white,minimum height=0.72cm,align=center,inner sep=4pt,line width=0.8pt},
  arr/.style={-{Latex[length=2mm]},thick}]
\node[n] (src) {字符流\\Source Code};
\node[n,right=1.2cm of src] (lex) {词法分析\\Lexical Analyzer};
\node[hi,right=1.2cm of lex] (par) {语法分析\\Parser};
\node[n,right=1.2cm of par] (sem) {语义分析\\Semantic Analyzer};
\node[n,right=1.2cm of sem] (ir) {中间代码\\IR};
\draw[arr] (src) -- (lex);
\draw[arr] (lex) -- node[above,font=\footnotesize]{token stream} (par);
\draw[arr] (par) -- node[above,font=\footnotesize]{parse tree} (sem);
\draw[arr] (sem) -- (ir);
\end{tikzpicture}
\end{center}

词法分析器把源代码拆成 token 流，例如 \pg{7}：
\[
\code{while (y<z)\{ int x=a+b; y+=x; \}}
\]
会被词法分析为类似
\[
\code{(keyword,while), (id,y), (sym,<), (id,z), ...}
\]
这样的 token 序列。语法分析器接手后，不再关心字符级细节，而是检查这些 token 是否构成合法语句。

\subsection{语法分析的定义 \pg{9,10}}
\begin{defbox}[语法分析（Syntax Analysis / Parsing）]
语法分析是编译的第二阶段，也称为\term{解析}{Parsing}或\term{语法分析器}{Parser}。
它从词法分析器获得 token 名组成的串，并验证该 token 串是否能由源语言的文法生成。
如果合法，它构造\term{分析树}{Parse Tree}并传给编译器前端后续阶段；如果非法，它报告语法错误。
\end{defbox}

更形式化地说，Parser 要完成两件事：
\begin{enumerate}
  \item \textbf{判定成员关系}：输入 token 串 \(w\) 是否属于某个文法 \(G\) 生成的语言 \(L(G)\)。
  \item \textbf{构造结构表示}：如果 \(w\in L(G)\)，构造它的推导结构，通常是 parse tree 或 AST。
\end{enumerate}

\subsection{语法分析举例 \pg{11}}
例 1：源程序输入：
\[
\code{if (x == y) stmt1 else stmt2}
\]
词法分析输出类似：
\[
\code{KEY(IF) '(' ID(x) OP('==') ID(y) ')' ... KEY(ELSE) ...}
\]
Parser 要判断它是否符合 \code{if} 语句的语法，并生成表示条件、then 分支、else 分支的树形结构。

例 2：token 串
\[
\code{<id,x> <op,*> <op,\%>}
\]
在 C 语言中不是合法语句。虽然每个 token 单独都可能合法，但它们的组合 \code{x *\%} 不符合表达式语法。

\begin{warnbox}[关键区分]
词法分析回答“每一段字符是什么 token”；语法分析回答“这些 token 能不能按语言规则组成句子”。合法 token 的任意拼接不一定是合法程序。
\end{warnbox}

% ============================================================
\section{如何定义语法：形式语言与文法（第 12--16 页）}
% ============================================================

\subsection{自然语言与形式语言 \pg{12,13}}
\begin{itemize}
  \item \term{自然语言}{Natural Language}：人类日常使用的语言，例如中文、英语。自然语言常有歧义，规则也可能依赖语境。
  \item \term{形式语言}{Formal Language}：由精确数学规则或机器可处理公式定义的语言，具有严格语法规则。
  \item \term{程序设计语言}{Programming Language}：一种形式语言，用于定义计算机程序。
\end{itemize}

形式语言可以用多种方式定义，例如\term{正则表达式}{Regular Expression}、
\term{有限自动机}{Finite Automata}和\term{文法}{Grammar}。其中 RE/FA 不够强，无法完整描述程序语法；文法尤其适合描述嵌套结构。

\subsection{文法的四元组定义 \pg{14}}
\begin{defbox}[文法（Grammar）]
文法是定义语言的数学模型。一个文法通常写作
\[
G=(V_T,V_N,S,\delta)
\]
其中 \(V_T\) 是终结符集合，\(V_N\) 是非终结符集合，\(S\) 是开始符号，\(\delta\) 是产生式集合，并且 \(V_T\cap V_N=\emptyset\)。
\end{defbox}

\begin{longtable}{>{\bfseries}p{0.18\linewidth}p{0.25\linewidth}p{0.48\linewidth}}
\toprule
组成 & 英文术语 & 含义 \\
\midrule
\(V_T\) & Terminals & \term{终结符}{terminal symbols}，构成最终字符串的基本符号。对编译器而言，终结符本质上通常就是 token，位于 parse tree 的叶子。\\
\(V_N\) & Non-terminals & \term{非终结符}{non-terminal symbols}，表示一类可继续展开的语法结构，如 expression、statement、loop，通常位于 parse tree 的内部节点。\\
\(S\) & Start Symbol & \term{开始符号}{start symbol}，推导的根，是某个非终结符。课件要求 \(S\) 至少出现在某条产生式左部一次。\\
\(\delta\) & Productions & \term{产生式}{productions}集合，规定终结符和非终结符如何组合成句子。\\
\bottomrule
\end{longtable}

\subsection{产生式与 BNF \pg{15}}
\begin{defbox}[产生式（Production）]
产生式形如
\[
\text{LHS}\to\text{RHS}
\]
读作“左部产生右部”或“左部定义为右部”。左部称为\term{产生式头}{head / left-hand side}，
右部称为\term{产生式体}{body / right-hand side}。\(\to\) 有时也写作 \(\mathrel{::=}\)，这就是常见的 \term{巴科斯范式}{Backus--Naur Form, BNF}记法。
\end{defbox}

产生式的直观含义是：当推导过程中遇到左部符号时，可以把它替换为右部符号串。

\subsection{中文句子的文法例子 \pg{16}}
课件给出一个简化中文句子文法：
\[
\begin{aligned}
\langle\text{句子}\rangle &\to \langle\text{主语}\rangle\langle\text{谓语}\rangle\langle\text{宾语}\rangle\\
\langle\text{主语}\rangle &\to \text{我}\mid\text{猫}\\
\langle\text{谓语}\rangle &\to \text{喜欢}\mid\text{追}\\
\langle\text{宾语}\rangle &\to \text{巧克力}\mid\text{老鼠}
\end{aligned}
\]

其中：
\[
V_N=\{\langle\text{句子}\rangle,\langle\text{主语}\rangle,\langle\text{谓语}\rangle,\langle\text{宾语}\rangle\}
\]
\[
V_T=\{\text{我},\text{猫},\text{喜欢},\text{追},\text{巧克力},\text{老鼠}\},\quad S=\langle\text{句子}\rangle
\]

它可以生成“猫喜欢巧克力”“我追老鼠”等句子。注意“老鼠追猫”若要被生成，则“老鼠”还必须能作为主语、“猫”还必须能作为宾语；这说明文法必须\textbf{精确且完整}，否则会漏掉本应合法的句子或生成不应合法的句子。

% ============================================================
\section{上下文无关文法 CFG（第 17--23 页）}
% ============================================================

\subsection{为什么 CFG 适合程序语法 \pg{17}}
检查程序是否\term{良构}{well-formed}，需要对合法程序结构给出规范。这个规范应满足：
\begin{itemize}
  \item \textbf{精确}：不能靠自然语言猜测。
  \item \textbf{完备}：覆盖语言所有语法细节。
  \item \textbf{便于使用}：语言设计者和编译器实现者都能使用。
\end{itemize}

\term{上下文无关文法}{Context-Free Grammar, CFG}正好满足这些需求。它足够描述当代大多数编程语言的主要语法结构，尤其是表达式、语句、块、嵌套结构等递归语法。

\subsection{CFG 的形式化定义 \pg{18}}
\begin{defbox}[上下文无关文法（Context-Free Grammar, CFG）]
CFG 仍写作四元组 \(G=(V_T,V_N,S,\delta)\)，但产生式必须形如
\[
A\to \alpha
\]
其中 \(A\in V_N\) 是\textbf{单个非终结符}，\(\alpha\in(V_N\cup V_T)^*\) 是终结符和非终结符组成的串。
\end{defbox}

“上下文无关”的意思是：只要看到非终结符 \(A\)，就可以把它替换成 \(\alpha\)，不需要查看 \(A\) 左右是什么上下文。

\subsection{算术表达式 CFG \pg{18,19,23}}
课件先给出一个只含 \(+\)、\(*\)、括号和标识符 \(i\) 的表达式文法：
\[
G=\langle \{i,+,*,(,)\},\{E\},E,\delta\rangle
\]
\[
E\to i \mid E+E \mid E*E \mid (E)
\]
这是一个能描述基本表达式的 CFG，但它没有规定 \(+\) 和 \(*\) 的优先级，也没有规定结合性，因此后面会产生二义性。

课件后面给出更常用的表达式文法：
\[
\begin{aligned}
E &\to E+T \mid E-T \mid T\\
T &\to T*F \mid T/F \mid F\\
F &\to (E)\mid id
\end{aligned}
\]
这里的 \(E,T,F\) 分别表示 expression、term、factor。它天然表达了“乘除优先于加减”：因为 \(+\) 和 \(-\) 只在 \(E\) 层出现，\(*\) 和 \(/\) 在更低的 \(T\) 层出现。

\subsection{文法书写约定 \pg{20--23}}
\begin{longtable}{p{0.32\linewidth}p{0.58\linewidth}}
\toprule
\textbf{符号} & \textbf{常见含义} \\
\midrule
\(a,b,c\) 等靠前小写字母 & 终结符。\\
\(+\), \(*\), \((\), \()\), 标点、数字 & 终结符。\\
\(\mathbf{id}\), \(\mathbf{if}\) 等粗体字符串 & 一个整体终结符，通常对应一个 token。\\
\(A,B,C\) 等靠前大写字母 & 非终结符。\\
\(S\) & 通常表示开始符号。\\
\(expr,stmt\) 等小写斜体名字 & 非终结符，表示程序构造。\\
\(E,T,F\) & 常用于表达式文法：表达式、项、因子。\\
\(X,Y,Z\) 等靠后大写字母 & 文法符号，可以是终结符或非终结符。\\
\(u,v,\ldots,z\) 等靠后小写字母 & 可能为空的终结符串。\\
\(\alpha,\beta,\gamma\) 等希腊字母 & 可能为空的文法符号串。\\
\bottomrule
\end{longtable}

共享同一左部的规则可以合并：
\[
\alpha\to\beta_1,\quad \alpha\to\beta_2,\quad\ldots,\quad\alpha\to\beta_n
\]
可写作：
\[
\alpha\to\beta_1\mid\beta_2\mid\cdots\mid\beta_n
\]
其中每个 \(\beta_i\) 称为一个\term{候选式}{alternative}。

默认情况下，如果未特别说明，第一条产生式的左部就是开始符号。

% ============================================================
\section{推导：从文法到语言（第 24--30 页）}
% ============================================================

\subsection{推导的基本定义 \pg{25}}
\begin{defbox}[推导（Derivation）]
推导是连续应用产生式规则的过程。若存在产生式 \(A\to\gamma\)，并且某个文法符号串中出现了 \(A\)，则
\[
\alpha A\beta \Rightarrow \alpha\gamma\beta
\]
表示“一步推导”，即把中间的 \(A\) 替换成 \(\gamma\)。
\end{defbox}

常用记号：
\begin{itemize}
  \item \(\alpha\Rightarrow\beta\)：一步推导。
  \item \(\alpha\Rightarrow^*\beta\)：零步或多步推导；因此任意 \(\alpha\Rightarrow^*\alpha\) 成立。
  \item \(\alpha\Rightarrow^+\beta\)：一步或多步推导。
  \item 若 \(\alpha\Rightarrow^*\beta\) 且 \(\beta\Rightarrow^*\gamma\)，则 \(\alpha\Rightarrow^*\gamma\)。
\end{itemize}

\subsection{句型、句子与语言 \pg{26}}
\begin{defbox}[句型、句子、语言]
设 \(S\) 是文法 \(G\) 的开始符号。
\begin{itemize}
  \item 若 \(S\Rightarrow^*\alpha\)，则 \(\alpha\) 是 \(G\) 的\term{句型}{sentential form}。句型可以包含终结符和非终结符，也可以为空。
  \item 若某个句型不含任何非终结符，则称为\term{句子}{sentence}。
  \item 文法生成的所有句子的集合称为该文法生成的\term{语言}{language}：
  \[
  L(G)=\{w\mid S\Rightarrow^*w,\; w\in V_T^*\}
  \]
\end{itemize}
\end{defbox}

\subsection[例子：A -> c | Ab 生成的语言]{例子：\(A\to c\mid Ab\) 生成的语言 \pg{27}}
文法：
\[
G[A]:\quad A\to c\mid Ab
\]
可以推导：
\[
A\Rightarrow c
\]
\[
A\Rightarrow Ab\Rightarrow cb
\]
\[
A\Rightarrow Ab\Rightarrow Abb\Rightarrow cbb
\]
\[
A\Rightarrow Ab\Rightarrow Abb\Rightarrow Abbb\Rightarrow cbbb
\]
因此它生成的语言是：
\[
L(G)=\{cb^n\mid n\ge 0\}
\]
也就是一个 \(c\) 后面跟任意多个 \(b\)。

\subsection{推导过程中的选择 \pg{28,29}}
每一步推导通常有两个选择：
\begin{enumerate}
  \item \textbf{替换哪个非终结符}。
  \item \textbf{对这个非终结符使用哪条产生式}。
\end{enumerate}

这两个选择直接对应 parser 的两大思路：
\begin{itemize}
  \item \term{自顶向下分析}{Top-down parsing}：从开始符号出发，尝试生成输入串。
  \item \term{自底向上分析}{Bottom-up parsing}：从输入串出发，尝试规约回开始符号。
\end{itemize}

\subsection{最左推导与最右推导 \pg{29,30}}
\begin{defbox}[最左推导与最右推导]
\term{最左推导}{Leftmost Derivation}：每一步总是替换当前句型中最左边的非终结符。\\
\term{最右推导}{Rightmost Derivation}：每一步总是替换当前句型中最右边的非终结符。
\end{defbox}

以课件文法
\[
E\to T\mid E+T,\quad T\to F\mid T*F,\quad F\to (E)\mid i
\]
推导 \((i*i+i)\) 为例。

\textbf{最左推导}：
\[
\begin{aligned}
E&\Rightarrow T\Rightarrow F\Rightarrow (E)\Rightarrow(E+T)\Rightarrow(T+T)\\
&\Rightarrow(T*F+T)\Rightarrow(F*F+T)\Rightarrow(i*F+T)\\
&\Rightarrow(i*i+T)\Rightarrow(i*i+F)\Rightarrow(i*i+i)
\end{aligned}
\]

\textbf{最右推导}：
\[
\begin{aligned}
E&\Rightarrow T\Rightarrow F\Rightarrow (E)\Rightarrow(E+T)\Rightarrow(E+F)\\
&\Rightarrow(E+i)\Rightarrow(T+i)\Rightarrow(T*F+i)\\
&\Rightarrow(T*i+i)\Rightarrow(F*i+i)\Rightarrow(i*i+i)
\end{aligned}
\]

两种推导替换顺序不同，但最终可对应同一棵分析树。

% ============================================================
\section{分析树 Parse Tree（第 31--35 页）}
% ============================================================

\subsection{分析树的定义 \pg{32,33}}
\begin{defbox}[分析树（Parse Tree）]
分析树是推导的图形化表示。它保留“用了哪些产生式、形成了什么层次结构”，但过滤掉“推导过程中先替换哪个非终结符”的时间顺序。
\end{defbox}

分析树的结构规则：
\begin{itemize}
  \item 每个内部节点表示一次产生式应用。
  \item 内部节点标签是产生式左部的非终结符。
  \item 它的子节点从左到右对应产生式右部的符号。
  \item 根节点通常是开始符号。
\end{itemize}

例如 \(E\Rightarrow E+T\) 在树中表现为：父节点 \(E\)，三个子节点依次为 \(E,+,T\)。

\subsection{叶子、产出与边缘 \pg{34,35}}
\begin{defbox}[Yield / Frontier]
分析树从左到右读出所有叶子节点，得到的符号串称为该树的\term{产出}{yield}或\term{边缘}{frontier}。它构成一个句型；如果全是终结符，就是一个句子。
\end{defbox}

\begin{notebox}[分析树与推导的关系]
普通推导可能有许多替换顺序，但分析树不记录这些顺序。对一棵固定分析树而言，可以唯一恢复出一条最左推导，也可以唯一恢复出一条最右推导。因此最左/最右推导常用于判定分析树是否唯一。
\end{notebox}

\subsection{小型分析树示意}
下面用一个简化树表示 \(E\to E+T\)、\(T\to F\)、\(F\to i\) 等规则如何组合：

\begin{center}
\begin{tikzpicture}[
  level distance=0.9cm,
  sibling distance=1.4cm,
  every node/.style={font=\small,align=center},
  edge from parent/.style={draw,-{Latex[length=1.4mm]}}
]
\node {E}
  child { node {E}
    child { node {T}
      child { node {F}
        child { node {i} }}}}
  child { node {+} }
  child { node {T}
    child { node {F}
      child { node {i} }}};
\end{tikzpicture}
\end{center}

它的叶子从左到右是 \(i+i\)，因此 yield 是 \(i+i\)。

% ============================================================
\section{二义性与消歧（第 36--41 页）}
% ============================================================

\subsection{二义性的基本问题 \pg{36,37}}
考虑文法：
\[
G(E):\quad E\to i\mid E+E\mid E*E\mid(E)
\]
句子 \((i*i+i)\) 可以有两种结构：
\[
(i*i)+i
\qquad\text{或}\qquad
i*(i+i)
\]
这对应两棵不同的分析树，也意味着同一程序片段可能有两种语义。

\begin{warnbox}[为什么程序语言不喜欢二义性]
如果一个程序字符串有两棵不同分析树，它就可能有两种不同含义。编译器必须给出确定行为，因此大多数 parser 更希望文法无二义，或者至少通过额外规则把多余分析树排除掉。
\end{warnbox}

\subsection{文法二义性与语言二义性 \pg{38}}
\begin{defbox}[文法二义性（Ambiguity of Grammar）]
如果某个文法存在一个句子，该句子对应多于一棵不同的分析树，则该文法是二义的。
\end{defbox}

\begin{defbox}[语言二义性（Ambiguity of Language）]
如果一个语言不存在任何无二义文法能生成它，则这个语言是二义语言，也称\term{固有二义}{inherently ambiguous}语言。
\end{defbox}

注意：某个文法 \(G\) 有二义性，不代表 \(L(G)\) 这个语言本身固有二义。可能存在另一个文法 \(G'\)，满足 \(L(G)=L(G')\)，但 \(G'\) 无二义。

\subsection{二义性判定的困难 \pg{39}}
课件指出：判断任意文法是否二义是\term{不可判定问题}{undecidable problem}。也就是说，不存在一个算法，能在有限步骤内对所有输入文法都准确判断“是否二义”。

因此工程上通常采用：
\begin{itemize}
  \item 人工重写文法，尽量消除二义性。
  \item 使用二义文法，但配合\term{消歧规则}{disambiguating rules}丢弃不想要的分析树，例如 YACC 中常见的优先级、结合性声明。
  \item 使用若干充分条件判断某类文法无二义；但充分条件不是必要条件，不能覆盖所有无二义文法。
\end{itemize}

\subsection{用优先级和结合性消除表达式二义性 \pg{40,41}}
表达式文法的二义性通常来自两个问题：
\begin{itemize}
  \item \term{优先级}{precedence}：例如 \(*\) 是否比 \(+\) 先结合。
  \item \term{结合性}{associativity}：例如 \(a-b-c\) 是 \((a-b)-c\) 还是 \(a-(b-c)\)。
\end{itemize}

一种无二义写法是：
\[
\begin{aligned}
E &\to E+T \mid T\\
T &\to T*F \mid F\\
F &\to (E)\mid i
\end{aligned}
\]
解释：
\begin{itemize}
  \item \(E\) 层处理 \(+\)，\(T\) 层处理 \(*\)，所以 \(*\) 优先级高于 \(+\)。
  \item \(E\to E+T\) 和 \(T\to T*F\) 使用左递归，通常表达左结合。
  \item 括号在 \(F\to(E)\) 中出现，能显式改变默认优先级。
\end{itemize}

% ============================================================
\section{乔姆斯基文法体系（第 42--49 页）}
% ============================================================

\subsection{四类文法总览 \pg{42,43}}
Noam Chomsky 在 1956 年建立形式语言体系，将文法按产生式限制分为 0、1、2、3 四类。它们都可写作 \(G=(V_T,V_N,S,\delta)\)，差异在于产生式允许的形式不同。

\begin{center}
\begin{tikzpicture}[font=\small,
  box/.style={draw,rounded corners=3pt,minimum width=8.5cm,minimum height=0.7cm,align=center},
  arr/.style={-{Latex[length=2mm]},thick}]
\node[box,fill=softbg] (t0) {Type 0：无限制文法 Unrestricted Grammar（最强）};
\node[box,fill=white,below=0.25cm of t0] (t1) {Type 1：上下文有关文法 Context-Sensitive Grammar};
\node[box,fill=softbg,below=0.25cm of t1] (t2) {Type 2：上下文无关文法 Context-Free Grammar};
\node[box,fill=white,below=0.25cm of t2] (t3) {Type 3：正则文法 Regular Grammar（最弱）};
\draw[arr] ($(t3.east)+(0.2,0)$) -- node[right,font=\footnotesize]{包含关系} ($(t0.east)+(0.2,0)$);
\end{tikzpicture}
\end{center}

包含关系可理解为：
\[
\text{Regular} \subset \text{Context-Free} \subset \text{Context-Sensitive} \subset \text{Unrestricted}
\]

\subsection{Type 0：无限制文法 \pg{44}}
\begin{defbox}[0 型文法（Unrestricted Grammar / Phrase-Structure Grammar）]
若每条产生式 \(\alpha\to\beta\) 满足：
\[
\alpha\in(V_N\cup V_T)^*,\quad \alpha\ne\varepsilon,\quad
\beta\in(V_N\cup V_T)^*
\]
则该文法是 0 型文法。
\end{defbox}

特点：
\begin{itemize}
  \item 左部可以是复杂符号串，不要求单个非终结符。
  \item 右部可以变长、变短，甚至为空。
  \item 与\term{图灵机}{Turing Machine}等价，表达能力最强。
  \item 推导过程中字符串长度可以反复扩张和收缩，达到目标串前的推导步数没有固定上界。
\end{itemize}

课件例子：
\[
aA\to aBCd,\qquad aBcd\to aE,\qquad A\to\varepsilon
\]

\subsection{Type 1：上下文有关文法 \pg{45}}
\begin{defbox}[1 型文法（Context-Sensitive Grammar, CSG）]
若每条产生式 \(\alpha\to\beta\) 满足：
\[
\alpha\in(V_N\cup V_T)^*,\quad \alpha\ne\varepsilon,\quad
\beta\in(V_N\cup V_T)^*,\quad |\alpha|\le|\beta|
\]
并且通常不允许空产生式（除非是受限的开始符号特例），则为上下文有关文法。
\end{defbox}

典型形式：
\[
\alpha A\beta\to\alpha\gamma\beta
\]
表示非终结符 \(A\) 只有在左边是 \(\alpha\)、右边是 \(\beta\) 这个上下文中，才能被替换为 \(\gamma\)。

特点：
\begin{itemize}
  \item 若不考虑特殊空串规则，推导串长度只能不减，因此推导步数受目标串长度约束。
  \item 与\term{线性有界自动机}{Linear Bounded Automaton, LBA}相关。
  \item 表达能力比 CFG 强，但解析复杂且效率低。
\end{itemize}

\subsection{Type 2：上下文无关文法 \pg{46}}
\begin{defbox}[2 型文法（Context-Free Grammar, CFG）]
若每条产生式形如
\[
A\to\alpha,\qquad A\in V_N,\quad \alpha\in(V_N\cup V_T)^*
\]
即左部是单个非终结符，则为上下文无关文法。
\end{defbox}

特点：
\begin{itemize}
  \item 替换 \(A\) 时不看上下文，因此称“上下文无关”。
  \item 对应\term{非确定下推自动机}{Nondeterministic Pushdown Automaton, NDPDA}。
  \item 能描述递归和嵌套结构，是语法分析的核心工具。
\end{itemize}

例子：
\[
S\to aSb\mid ab
\]
生成
\[
L=\{a^n b^n\mid n\ge 1\}
\]
该语言不是正则语言，但可以由 CFG 描述。

\subsection{Type 3：正则文法 \pg{47,48}}
\begin{defbox}[3 型文法（Regular Grammar）]
若每条产生式形如
\[
A\to aB\quad\text{或}\quad A\to a
\]
其中 \(A,B\in V_N\)，\(a\in V_T\cup\{\varepsilon\}\)，则为正则文法。开始符号可在受限条件下推出 \(\varepsilon\)。
\end{defbox}

特点：
\begin{itemize}
  \item 左部是单个非终结符，右部最多是一个终结符加一个非终结符。
  \item 与\term{有限自动机}{Finite Automaton, FA}等价。
  \item 推导串长度通常每步增加 1，因此表达能力最弱。
\end{itemize}

课件例子：
\[
A\to 1A\mid 0
\]
对应正则表达式：
\[
1^*0
\]

\subsection{四类文法对比 \pg{49}}
\begin{longtable}{p{0.12\linewidth}p{0.2\linewidth}p{0.27\linewidth}p{0.28\linewidth}}
\toprule
类型 & 名称 & 产生式限制 & 对应模型 / 用途 \\
\midrule
0 型 & 无限制文法 & \(\alpha\to\beta\)，左部非空，几乎不限制长度变化 & 图灵机；理论表达能力最强。\\
1 型 & 上下文有关文法 & \(|\alpha|\le|\beta|\)，基本不收缩 & 线性有界自动机；表达上下文约束但解析效率差。\\
2 型 & 上下文无关文法 & \(A\to\alpha\)，左部单个非终结符 & 下推自动机；编译器语法分析核心。\\
3 型 & 正则文法 & \(A\to aB\) 或 \(A\to a\) & 有限自动机；词法分析核心。\\
\bottomrule
\end{longtable}

% ============================================================
\section{实践中的语法分析与 AST（第 50--55 页）}
% ============================================================

\subsection{为什么编译器仍大量使用 CFG \pg{50}}
课件指出，大多数现代程序设计语言整体上并不是纯上下文无关语言，甚至很多约束也不是上下文有关文法适合直接处理的。例如：
\begin{itemize}
  \item 变量必须先声明再使用。
  \item 函数调用参数数量和类型必须匹配。
  \item 标识符作用域、类型检查、继承关系等需要语义信息。
\end{itemize}

但 CFG 仍被广泛用于编译器：
\begin{itemize}
  \item CFG 非常适合描述表达式和语句的递归结构。
  \item CSG parser 理论和实现都复杂，效率不适合常规编译。
  \item CFG 的构造和解析技术成熟高效。
  \item 剩余的上下文相关约束通常放到\term{语义分析}{Semantic Analysis}阶段处理。
\end{itemize}

\begin{notebox}[工程分工]
在程序语言实现中，常见分工是：
\[
\text{正则语言} \to \text{词法分析}
\qquad
\text{上下文无关语言} \to \text{语法分析}
\qquad
\text{上下文/类型约束} \to \text{语义分析}
\]
\end{notebox}

\subsection{Parsing 到底做什么 \pg{51}}
语法分析可以概括为：基于给定文法处理输入串，如果输入串属于该文法生成的语言，就构造它的推导表示；否则报告错误。

两个子任务：
\begin{enumerate}
  \item 判断一个字符串是否能由文法推导出来。
  \item 构造推导的表示，并传给编译器下一阶段。
\end{enumerate}

推导表示可以是：
\begin{itemize}
  \item \term{分析树}{Parse Tree}：完整保留文法推导细节。
  \item \term{抽象语法树}{Abstract Syntax Tree, AST}：省略不影响语义的细节，更适合后续编译阶段。
\end{itemize}

\subsection{Parse Tree vs AST \pg{52}}
\begin{defbox}[抽象语法树（Abstract Syntax Tree, AST）]
AST 是 parse tree 的缩略表示。它在不改变程序含义的前提下，删除括号、某些中间非终结符、纯语法辅助节点等推导细节，只保留语义上重要的结构。
\end{defbox}

\begin{longtable}{p{0.28\linewidth}p{0.3\linewidth}p{0.32\linewidth}}
\toprule
对比项 & Parse Tree & AST \\
\midrule
关注点 & 文法如何一步步推导出输入串 & 程序真正的语义结构 \\
节点数量 & 多，包含所有终结符、非终结符和标点 & 少，省略无意义中间节点 \\
括号与标点 & 通常保留 & 只在影响语义时保留 \\
适用阶段 & 语法分析解释、教学、调试 & 语义分析、中间代码生成、优化 \\
\bottomrule
\end{longtable}

例如表达式 \((i+i)+i\) 的 parse tree 会包含括号、\(E,T,F\) 等许多中间节点；AST 通常只保留两个 \(+\) 节点和三个 \(i\) 叶子：
\[
\begin{array}{c}
\texttt{+}\\[-2pt]
/ \quad \backslash\\[-2pt]
\texttt{+}\quad \texttt{i}\\[-2pt]
/ \backslash\\[-2pt]
\texttt{i}\quad \texttt{i}
\end{array}
\]

\subsection{课程总结 \pg{53--55}}
本讲的核心知识点可以压缩为：
\begin{itemize}
  \item 语法分析输入 token 流，判断它是否属于文法生成的语言。
  \item 文法 \(G=(V_T,V_N,S,\delta)\) 是定义形式语言的数学工具。
  \item CFG 的产生式形如 \(A\to\alpha\)，适合描述程序语言的递归语法结构。
  \item 推导从开始符号出发，连续应用产生式生成句型和句子。
  \item Parse tree 展示字符串如何由文法推导而来。
  \item 文法二义性意味着某个句子有多棵 parse tree，通常要用优先级、结合性或文法重写消除。
  \item AST 是程序语法结构的抽象表示，比 parse tree 更适合后续编译阶段。
  \item 延伸阅读可参考 Dragon Book 第 2.4、4.1.1、4.2、4.3 节。
\end{itemize}

% ============================================================
\section{练习题与参考答案}
% ============================================================

\subsection*{练习 1：区分词法分析和语法分析}
源程序片段为：
\[
\code{x * \%}
\]
词法分析器可能把它识别为 \(\code{id(x)},\code{op(*)},\code{op(\%)}\)。请问为什么它仍然可能被 parser 判为非法？

\textbf{答案：}词法分析只判断每个字符片段能否构成合法 token；语法分析要判断 token 序列是否符合表达式或语句文法。\(\code{id}\ \code{*}\ \code{\%}\) 中 \(\code{*}\) 后面缺少合法操作数，\(\code{\%}\) 也不能直接作为操作数，因此 token 各自合法，但组合不合法。

\subsection*{练习 2：写出文法四元组}
给定文法：
\[
E\to E+T\mid T,\quad T\to T*F\mid F,\quad F\to(E)\mid id
\]
写出它的 \(V_T,V_N,S\)。

\textbf{答案：}
\[
V_N=\{E,T,F\},\quad V_T=\{+,*,(,),id\},\quad S=E
\]
\(\delta\) 是题目给出的全部产生式集合。

\subsection*{练习 3：由文法推导字符串}
文法 \(G[A]:A\to c\mid Ab\)。推导字符串 \(cbbb\)。

\textbf{答案：}
\[
A\Rightarrow Ab\Rightarrow Abb\Rightarrow Abbb\Rightarrow cbbb
\]
因此 \(cbbb\in L(G)\)。

\subsection*{练习 4：判断句型和句子}
设 \(S\Rightarrow^* aAb\Rightarrow aabb\)，其中 \(A\) 是非终结符，\(a,b\) 是终结符。问 \(aAb\) 和 \(aabb\) 分别是什么？

\textbf{答案：}
\(aAb\) 是句型，因为它可由开始符号推导出，但仍含非终结符 \(A\)。\(aabb\) 是句子，因为它只含终结符。

\subsection*{练习 5：最左推导}
使用文法
\[
E\to T\mid E+T,\quad T\to F\mid T*F,\quad F\to(E)\mid i
\]
给出 \((i*i+i)\) 的最左推导。

\textbf{答案：}
\[
\begin{aligned}
E&\Rightarrow T\Rightarrow F\Rightarrow (E)\Rightarrow(E+T)\Rightarrow(T+T)\\
&\Rightarrow(T*F+T)\Rightarrow(F*F+T)\Rightarrow(i*F+T)\\
&\Rightarrow(i*i+T)\Rightarrow(i*i+F)\Rightarrow(i*i+i)
\end{aligned}
\]

\subsection*{练习 6：解释二义性}
为什么文法
\[
E\to i\mid E+E\mid E*E\mid(E)
\]
对表达式 \((i*i+i)\) 是二义的？

\textbf{答案：}
因为该文法没有规定 \(*\) 和 \(+\) 的优先级，也没有规定结合性。\((i*i+i)\) 可以被解释为 \((i*i)+i\)，也可以被解释为 \(i*(i+i)\)，对应两棵不同 parse tree。因此文法对该句子二义。

\subsection*{练习 7：改写无二义表达式文法}
写出一个能体现 \(*\) 优先级高于 \(+\)，并且二者左结合的表达式文法。

\textbf{答案：}
\[
E\to E+T\mid T,\quad T\to T*F\mid F,\quad F\to(E)\mid i
\]
\(E\) 层处理加法，\(T\) 层处理乘法，因此乘法优先级更高；左递归结构体现左结合。

\subsection*{练习 8：判断文法类型}
判断以下文法属于哪一类乔姆斯基文法：
\[
S\to aSb\mid ab
\]

\textbf{答案：}
它是 2 型文法，即上下文无关文法。每条产生式左部都是单个非终结符 \(S\)，右部是终结符和非终结符组成的串。它生成 \(L=\{a^n b^n\mid n\ge1\}\)，该语言不是正则语言。

\subsection*{练习 9：为什么不用 CSG 做主要语法分析}
既然上下文有关文法比 CFG 更强，为什么编译器主要使用 CFG？

\textbf{答案：}
CSG 表达能力更强，但解析复杂、效率低，不适合作为常规编译器主力工具。CFG 已经足够描述表达式、语句、块结构等主要递归语法，而且解析技术成熟高效。变量声明、类型匹配、作用域等上下文相关约束通常交给语义分析阶段处理。

\subsection*{练习 10：Parse Tree 与 AST}
为什么编译器后续阶段通常更喜欢 AST，而不是完整 parse tree？

\textbf{答案：}
Parse tree 完整展示文法推导，包含许多中间非终结符、括号、标点等语法细节。AST 删除不影响语义的细节，只保留程序结构和语义关系，因此更紧凑、更适合语义分析、中间代码生成和优化。

\begin{notebox}[考前速记]
\begin{itemize}
  \item Parser 输入 token 流，输出 parse tree 或 AST。
  \item Grammar：\(G=(V_T,V_N,S,\delta)\)。
  \item CFG：\(A\to\alpha\)，左部是单个非终结符。
  \item \(S\Rightarrow^*\alpha\)：\(\alpha\) 是句型；若 \(\alpha\in V_T^*\)，则是句子。
  \item Parse tree 的叶子从左到右组成 yield/frontier。
  \item 二义文法：某个句子有多棵 parse tree。
  \item 表达式消歧：靠优先级、结合性，或重写为 \(E/T/F\) 分层文法。
  \item Chomsky：正则 \(\subset\) 上下文无关 \(\subset\) 上下文有关 \(\subset\) 无限制。
\end{itemize}
\end{notebox}

\end{document}
